\section{Introduction}

For Riemannian symmetric spaces $G/K$ of the compact or non-compact type, there is a~well-known contraction principle
which states that under suitable scaling, the spherical functions $\varphi_\lambda$ of $G/K$ tend to the spherical
functions $\psi_\lambda$ of the tangent space of $G/K$ in the base point, which is a~symmetric space of the f\/lat type:
\begin{gather*}
\lim_{n\to\infty} \varphi_{n\lambda}(\exp(x/n))=\psi_\lambda(x).
\end{gather*}
See~\cite{C} and, for a~more recent account,~\cite{BO}.
This curvature limit, also known as Mehler--Heine formula, extends to the more general setting of hypergeometric
functions associated with root systems, which converge under rescaling to generalized Bessel functions.
This is proven in~\cite{dJ} by a~limit transition in the Cherednik operators; see also~\cite{BO} for a~dif\/ferent
approach.
In the compact rank one case, the contraction principle is a~weak version of the classical Hilb formula for Jacobi
polynomials (see~\cite[Theorem 8.21.12]{Sz}), which provides in addition a~precise estimate on the rate of
convergence.
In this paper, we prove a~Mehler--Heine formula with a~precise estimate on the error term
for a~certain class of orthogonal polynomials associated with root systems, which in particular encompasses the
spherical functions of compact Grassmannians.
This result is a~``compact'' analogue of Theorem 5.4 in~\cite{RV2}, which gives a~scaling limit with error bounds for
hypergeometric functions in the dual, non-compact setting.
In the second part of the paper, we shall use the Mehler--Heine formula~\ref{vgl-Bessel} in order to establish a~central
limit theorem for random walks on the semi-lattice of dominant weights parametrizing the unitary dual of a~compact
Grassmannian.

To become more precise, we consider the compact Grassmannians $\mathcal G_{p,q}(\mathbb F)=G/K$ over one of the (skew-)
f\/ields $\mathbb F= \mathbb R, \mathbb C, \mathbb H$, with $G=SU(p+q,\mathbb F)$ and $K=S(U(q, \mathbb F)\times
U(p,\mathbb F))$, where $p\ge q\ge 1$.
Via polar decomposition of~$G$, the double coset space $G//K=\{KgK:\> g\in G\}$ may be topologically identif\/ied with the
fundamental alcove
\begin{gather*}
A_0:=\Big\{x=(x_1,\ldots,x_q)\in\mathbb R^q:\> \frac{\pi}{2}\ge x_1\ge x_2\ge\dots\ge x_q\ge0\Big\},
\end{gather*}
with $x\in A_0$ being identif\/ied with the matrix
\begin{gather*}
a_x=
\begin{pmatrix}
\cos\underline x & -\sin \underline x & 0
\\
\sin \underline x & \cos \underline x & 0
\\
0 & 0 & I_{p-q}
\end{pmatrix}
.
\end{gather*}
Here we use the diagonal matrix notation $\underline x=\diag(x_1,\ldots,x_q)$, and the functions $\sin$, $\cos$ are
understood component-wise.
For details, see~\cite[Theorem~4.1]{RR}.
The spherical functions of~$\mathcal G_{p,q}(\mathbb F)$ can be viewed as Heckman--Opdam polynomials of type~$BC_q$,
which are also known as multivariable Jacobi polynomials.
They may be described as follows: denote by~$F_{BC}(\lambda,k;\cdot)$ the Heckman--Opdam hypergeometric function
associated with the root system
\begin{gather*}
R=2  BC_q=\{\pm 2e_i, \pm 4 e_i: \> 1\leq i\le q\}\cup\{\pm 2e_i \pm 2e_j: 1\leq i < j \leq q\} \subset \mathbb R^q,
\end{gather*}
with spectral variable $\lambda\in\mathbb C^q$ and multiplicity parameter $k=(k_1, k_2, k_3)\in \mathbb R^3$
corresponding to the roots $\pm 2 e_i$, $\pm 4 e_i$ and $2(\pm e_i\pm e_j)$.
Fix the positive subsystem
\begin{gather*}
R_+=\{2e_i, 4e_i, 1\leq i \leq q\}\cup\{2e_i\pm 2e_j, 1\leq i<j \leq q\}
\end{gather*}
and the associated semi-lattice of dominant weights,
\begin{gather*}
P_+=\big\{\lambda\in (2\mathbb Z)^q: \> \lambda_1\ge\lambda_2\ge\dots\ge\lambda_q\ge0\big\}.
\end{gather*}
Then the set of spherical functions of $\mathcal G_{p,q}(\mathbb F)$ is parametrized by $P_+$ and consists of the
functions
\begin{gather}
\label{jacobi-pol-def}
\varphi_\lambda^p(a_x)=F_{BC}(\lambda+\rho_p, k(p),ix)=: R_\lambda^p(x),
\qquad
\lambda\in P_+
\end{gather}
with multiplicity parameter
\begin{gather}
\label{multipl}
k(p)=(d(p-q)/2, (d-1)/2, d/2),
\end{gather}
where $d=\dim_\mathbb R \mathbb F \in \{1,2,4\}$ and
\begin{gather*}
\rho_p=\frac{1}{2}\sum\limits_{\alpha\in R_+} k_\alpha\alpha=\sum\limits_{i=1}^q \left(\frac{d}{2}(p+q+2-2i) -1\right)e_i.
\end{gather*}
The functions $R_\lambda^p$ are the Heckman--Opdam polynomials associated with the root system~$R$ (called Jacobi
polynomials in the following) and with multiplicity $k(p)$, normalized according to $R_\lambda^p(0)=1$.
We refer to~\cite{H,HS,O1} for Heckman--Opdam theory in general, and to~\cite{RR} and the references
cited there for the connection with spherical functions in the compact $BC$ case.
Notice that our notion of~$F_{BC}$ coincides with that of Heckman, Opdam and~\cite{R2,RV2}, while it dif\/fers
from the geometric notion in~\cite{RR}.
Theorem~4.3 of~\cite{RR} corresponds to~\eqref{jacobi-pol-def}.

In Theorem 4.2 of~\cite{RR}, the product formula for spherical functions of $(G,K)$ was written as a~formula on~$A_0$
and analytically extended to a~product formula for the Jacobi polyno\-mials~$R_\lambda^p$ with multiplicity~$k(p)$
corresponding to arbitrary real parameters $p>2q-1$.
This led to three continuous series of positive product formulas for Jacobi polynomials corresponding to $\mathbb
F=\mathbb R, \mathbb C, \mathbb H$ and to associated commutative hypergroup structures on $A_0$; see~\cite{J}
and~\cite{BH} for the notion of hypergroups.
Using a~Harish-Chandra-type integral representation for the $R_\lambda^p$, we shall derive a~Mehler--Heine formula with
a~precise asymptotic estimate for the Jacobi polyno\-mials~$R_{\lambda}^p$ in terms of Bessel functions associated with
root system $B_q$ on the Weyl chamber
\begin{gather*}
C=\{x= (x_1, \ldots, x_q) \in \mathbb R^q: x_1 \geq \cdots \geq x_q \geq 0\}.
\end{gather*}
This Mehler--Heine formula will be the key ingredient for the main result of the present paper, a~central limit theorem
for random walks on the semi-lattice $P_+$, which parametrizes the spherical unitary dual of $G/K$.
To explain this CLT, let us f\/irst recall that via the GNS representation, the spherical functions $\varphi_\lambda,
\lambda \in P_+$ of $(G,K)$, which are necessarily positive def\/inite, are in a~one-to-one correspondence with the
(equivalence classes of) spherical representations $(\pi_\lambda,H_\lambda)$ of~$G$, that is those irreducible unitary
representations of~$G$ whose restriction to~$K$ contains the trivial representation with multiplicity one, see~\cite{F}
or~\cite[Chapter~IV]{Hel}.
The decomposition of tensor products of spherical representations into their irreducible components leads to
a~probability preserving convolution~$*_{d,p}$ and f\/inally a~Hermitian hypergroup structure on the discrete set~$P_+$;
see~\cite{Du} and~\cite{K3}.
Following, e.g.,~\cite{BH,V1,Z}, we introduce random walks $(S_n^{d,p})_{n\ge0}$ on $P_+$ associated
with $*_{d,p}$ and derive some limit theorems for $n\to\infty$.
The main result of this paper will be the Central Limit Theorem~\ref{clt}.
